% Example 8.2: the same weld length arranged three ways, seen looking along the bracket.
\begin{tikzpicture}[>=Latex,line join=round]
\tikzset{lb/.style={font=\normalsize}, sm/.style={font=\small}, wd/.style={line width=2.8pt,line cap=round}, sec/.style={fill=black!10,draw=black,line width=0.5pt}}
% side view
\begin{scope}
 \fill[black!25] (-0.7,-1.9) rectangle (0,1.9); \draw[line width=0.6pt] (-0.7,-1.9) rectangle (0,1.9);
 \fill[black!10] (0,-1.0) rectangle (3.4,1.0); \draw[line width=0.6pt] (0,-1.0) rectangle (3.4,1.0);
 \draw[wd] (0,1.0) -- (0,-1.0);
 \draw[->,line width=1.4pt] (3.0,1.9) -- (3.0,1.05); \node[lb,anchor=south] at (3.0,1.9) {$P=\SI{8}{kN}$};
 \draw[<->,line width=0.4pt] (0,-1.4) -- (3.0,-1.4); \node[sm,anchor=north] at (1.5,-1.42) {150};
 \draw[<->,line width=0.4pt] (3.7,-1.0) -- (3.7,1.0); \node[sm,anchor=west] at (3.75,0) {100};
 \node[lb,anchor=north] at (1.4,-2.2) {side view};
\end{scope}
% three weld patterns, end view
\foreach \xo/\ttl/\iu/\tv in {7.0/{two vertical}/{$d^3/6$}/85, 10.6/{two horizontal}/{$bd^2/2$}/30, 14.2/{all round}/{$d^2(3b+d)/6$}/22}{
 \begin{scope}[xshift=\xo cm]
  \path[sec] (-1.0,-1.0) rectangle (1.0,1.0);
  \draw[line width=0.3pt,dash pattern=on 5pt off 2pt on 1pt off 2pt] (-1.5,0) -- (1.5,0);
  \node[lb,anchor=north] at (0,-1.3) {\ttl};
  \node[sm,anchor=north] at (0,-1.95) {$I_u=$ \iu};
  \node[lb,anchor=north] at (0,-2.65) {$\tau=\SI{\tv}{MPa}$};
 \end{scope}}
\begin{scope}[xshift=7.0cm] \draw[wd] (-1.0,-1.0) -- (-1.0,1.0) (1.0,-1.0) -- (1.0,1.0); \end{scope}
\begin{scope}[xshift=10.6cm] \draw[wd] (-1.0,1.0) -- (1.0,1.0) (-1.0,-1.0) -- (1.0,-1.0); \end{scope}
\begin{scope}[xshift=14.2cm] \draw[wd] (-1.0,-1.0) rectangle (1.0,1.0); \end{scope}
\node[sm,anchor=west] at (15.75,0) {neutral axis};
\end{tikzpicture}
